\begin{trapbox}
	\begin{description}[style=nextline, leftmargin=2.8em, labelsep=0.5em, itemsep=0.35em]
		\item[\textbf{\textcolor{CTrap}{易错点一（忽略边界条件）}}] 认为递推公式只对 $n\geq m\geq 1$ 成立，忽略 $m=0$ 或 $m=n+1$ 时的边界情况。
		\item[\textbf{\textcolor{CTrap}{正确理解（约定补充边界）：}}] 在约定 $\binom{n}{k}=0$（$k<0$ 或 $k>n$）下，公式对一切非负整数 $n,m$ 自动成立。
		\item[\textbf{\textcolor{CTrap}{错因分析（机械记忆条件）：}}] 机械记忆条件，未理解约定对公式完整性的作用。
	\end{description}

	\medskip
	\begin{description}[style=nextline, leftmargin=2.8em, labelsep=0.5em, itemsep=0.35em]
		\item[\textbf{\textcolor{CTrap}{易错点二（混淆递推方向）}}] 将公式误用为 $\binom{n+1}{m} = \binom{n}{m} + \binom{n+1}{m+1}$ 等错误方向。
		\item[\textbf{\textcolor{CTrap}{正确理解（分清左右两边）：}}] 递推公式是左边两项和等于右边一项，关系为 $\binom{n}{m}+\binom{n}{m-1}=\binom{n+1}{m}$，顺序不能颠倒。
		\item[\textbf{\textcolor{CTrap}{错因分析（下标规律不清）：}}] 未看清下标增减规律，混淆了加法与裂项的方向。
	\end{description}

	\medskip
	\begin{description}[style=nextline, leftmargin=2.8em, labelsep=0.5em, itemsep=0.35em]
		\item[\textbf{\textcolor{CTrap}{易错点三（忽略定义域）}}] 直接使用公式时未检查 $n,m$ 是否满足 $n\geq m$，导致出现负数下标或无意义组合数。
		\item[\textbf{\textcolor{CTrap}{正确理解（确认变量范围）：}}] 使用前确认 $n,m$ 为非负整数且 $n\geq m$；若出现越界应利用约定处理（该补零）。
		\item[\textbf{\textcolor{CTrap}{错因分析（定义域不敏感）：}}] 对组合数定义域不敏感，忽略公式成立的前提。
	\end{description}
\end{trapbox}
